\section{Large Multi-Modal Models}

\textbf{Large Multi-Modal Models} (LMMs) are a class of models that can process and generate data across multiple modalities, such as text, images, audio, and video. 
In particular, we will focus on \textbf{Large Vision-Language Models} (LVLMs), which are models that can understand and generate both visual and textual information, 
usually by combining a \textbf{vision encoder} (converting images into embeddings) and a \textbf{Large Language Model}, to address various tasks.

\begin{comment}
    # TODO: add a figure showing the architecture of a LVLM, with a vision encoder and a LLM, and how they interact to process multi-modal data.
\end{comment}

We will first discuss different approaches to designing and training \textbf{vision encoder} architectures for LVLMs.
Finally, we will discuss the architectures that allow for the integration of vision and language modalities, and how to train these models to perform various tasks.

\subsection{Encoders}
Over the years, several architectures have been proposed for vision encoders. We can group them into four main categories: 
\begin{itemize}
    \item \textbf{Convolutional Neural Networks (CNNs)}: based on supervised learning, these models are limited by the fact that the labels need to be pre-defined, and they are not able to learn from unlabeled data.
    \item \textbf{Image-Only (Non-)Contrastive Learning}: based on self-supervised learning, these models learn to generate embeddings for images without the need for labels.
    This is achieved by training the model to output similar embeddings for different augmentations of the same image.
    \item \textbf{Masked Image Modeling (MIM)}: given an image, the model is trained to predict the missing parts of the image, which allows it to learn a representation of the image that captures its structure and semantics.
    \item \textbf{Contrastive Learning Approaches}: architectures like \textbf{CLIP}, that learn to align image and text embeddings in a shared space, allowing for cross-modal retrieval and understanding.
\end{itemize}

We will discuss each of these approaches in more detail in the following sections.

\subsection{Image-Only (Non-)Contrastive Learning}
Various self-supervised learning approaches have been proposed to learn image representations without the need for labels.

\subsubsection{SimCLR}
SimCLR is a contrastive learning approach that learns to generate embeddings for images by training the model to output similar embeddings for different augmentations of the same image, while ensuring that embeddings for different images are dissimilar.
Given an image, the model first generates two different augmentations of the image. 
A base encoder is then used to generate embeddings for both augmentations. 
These embeddings are then passed through a projection head, which is trained to maximize the similarity between the embeddings of the two augmentations, while minimizing the similarity between embeddings of different images.

\begin{comment}
    #TODO: add a figure showing the architecture of SimCLR, with the base encoder, projection head, and the contrastive loss function.
\end{comment}

For downstream tasks, the projection head is discarded, and the base encoder is used to generate embeddings for images.

The used loss function is the \textbf{InfoNCE loss}, which is a contrastive loss that encourages the model to output similar embeddings for positive pairs (different augmentations of the same image) and dissimilar embeddings for negative pairs (different images).
It is defined as follows:
\begin{equation}
    \mathcal{L}_{\text{InfoNCE}} = -\log \frac{\exp(\text{sim}(z_i, z_j) / \tau)}{\sum_{k=1}^{N} \exp(\text{sim}(z_i, z_k) / \tau)}
\end{equation}
where we are basically maximizing the similarity between the embeddings of the two augmentations of the same image ($z_i$ and $z_j$), while minimizing the similarity between embeddings of different images ($z_i$ and $z_k$), with $\tau$ being a temperature parameter that controls the sharpness of the distribution.

Such approaches have a main limitation: a large batch size is required to have a sufficient number of negative pairs, which can be computationally expensive and memory-intensive.


\subsubsection{DINO}
Recent approaches such as \textbf{DINO} (DIstillation with NO labels) have been proposed to overcome some limitations of contrastive learning methods such as SimCLR. 
DINO is a self-supervised learning framework based on Vision Transformers (ViTs) that learns meaningful image representations without requiring labeled data or negative sample pairs.

The architecture of DINO consists of two networks: a \textbf{student network} and a \textbf{teacher network}. 
The student network is trained to generate embeddings from input images, while the teacher network produces target embeddings that the student aims to match. 
Rather than being updated through backpropagation, the teacher's parameters are updated as an exponential moving average (EMA) of the student network's parameters, 
providing stable and slowly evolving targets throughout training.

During training, multiple augmented views of the same image are generated. 
The student network processes both global and local image crops, whereas the teacher network processes only the global crops. 
Both networks output feature vectors, which are projected through a small multi-layer perceptron (MLP) head to produce logits. 
These logits are converted into probability distributions using a softmax function.

The student is trained to match the probability distribution produced by the teacher for different views of the same image. 
This is achieved by minimizing the cross-entropy loss between the teacher's output distribution $P_t$ and the student's output distribution $P_s$:

\[
\mathcal{L} = - \sum_{i} P_t(i)\log P_s(i),
\]

where $i$ indexes the output dimensions. Since the teacher's output serves as the target distribution, gradients are computed only with respect to the student network.

The resulting embeddings have been shown to transfer effectively to a wide range of downstream computer vision tasks, 
including image classification, object detection, and semantic segmentation.


\subsection{Masked Image Modelling}
This family of approaches is based on the idea of predicting missing parts of an image, which allows the model to learn a representation of the image that captures its structure and semantics.

\subsubsection{BEiT}
\textbf{BERT Pre-Training of Image Transformers} (BEiT) adapts the masked language modeling objective of BERT to computer vision. 
First, an image is divided into non-overlapping patches. 
A pretrained image tokenizer converts each patch into a discrete visual token, which serves as the prediction target during pre-training. 
Meanwhile, the image patches are embedded as in a standard Vision Transformer, and a subset of the patch embeddings is replaced with a learned mask embedding. 

The Vision Transformer processes the resulting sequence together with positional embeddings. 
During training, the model predicts the discrete visual tokens corresponding to the masked patches by minimizing the cross-entropy loss between the predicted token distributions and the tokenizer-generated targets.

\subsubsection{MAE}
\textbf{Masked Autoencoders} (MAE) is a self-supervised learning approach that learns to reconstruct missing parts of an image.
Instead of predicting discrete visual tokens, MAE directly reconstructs the pixel values of the masked patches.

Usually, a large portion of the input image patches (e.g., 75\%) is randomly masked, and the remaining visible patches are processed by an encoder.
A decoder then takes the encoded visible patches and the mask tokens (which are included after the encoder) to reconstruct the original image.

\subsection{CLIP}
We have already discussed the \textbf{CLIP} architecture in the previous chapter, which is a contrastive learning approach that learns to align image and text embeddings in a shared space.
Although it shows impressive performance on various tasks, there are still some remarkable limitations that need to be addressed:
\begin{itemize}
    \item \textbf{No direct interaction between modalities:} CLIP simply learns to align the embeddings of images and text.
    It doesn't allow to go from one modality to another, or to generate new data in a different modality.
    \item \textbf{Limited task performance:} CLIP can only perform tasks that can be solved by comparing embeddings, such as image classification or retrieval.
    \item \textbf{No generative capabilities:} CLIP lacks the ability to generate new data, making them unsuitable for tasks that require data generation, such as captioning or image synthesis.
\end{itemize}

In this section, we will discuss some approaches that address these limitations and allow for more direct interaction between modalities, as well as generative capabilities.

\subsubsection{Data Scaling}
Various studies have shown that increasing the amount of training data can significantly improve the performance of LVLMs.
For example, scaling from 400 million to 2 billion image-text pairs has been shown to improve the performance of CLIP on various tasks.

Scale is not the only factor that affects the performance of LVLMs. Also, the quality of the data is important. 
Smaller, more curated datasets were shown to outperform larger, noisier datasets.

\subsubsection{Model Design}
Not only the amount and quality of data, but also the design of the model architecture can significantly affect the performance of CLIP.
In particular, new architectures have been proposed for both the vision encoder and the language model:

\begin{itemize}
    \item \textbf{Image Side:} the vision encoder can be improved via the use \textbf{masking}. 
    It has been shown that masking a large portion of the input image and only encoding the remaining visible patches
    can improve the performance of the vision encoder, reducing also the traing time and computational requirements.

    \item \textbf{Text Side:} along with the original text input, additional information such as the definition of the class or a description of the image can be provided to the language model, 
    which can improve the performance of the model on various tasks. This was particularly effective with fine-grained datasets, which typically contain many rare classes.
\end{itemize}

Some efforts have also been made to improve the interaction between modalities, while maintaining the shared embedding space.

For example, the \textbf{ImageBind} architecture allows for the integration of multiple modalities, such as images, text, audio, and video, into a single model.
To achieve this, a pretrained version of CLIP is used and kept frozen. A new encoder is used for each new modality, 
and a contrastive loss is used to align the embeddings of the new modalities with the embeddings of the original CLIP model.

\subsubsection{Interpretability}
One of the main limitations of CLIP is the limited interpretability of its learned embeddings. 
To address this issue, \textbf{STAIR} introduces a sparse and semantically meaningful representation for both image and text embeddings.

STAIR constructs a concept-based embedding space in which each dimension corresponds to a human-interpretable semantic concept. 
Instead of representing an image or a text as a dense feature vector, each embedding is expressed as a non-negative set of concept activations, 
where the value associated with each concept reflects its relevance to the input.

To obtain this representation, STAIR employs a \textbf{token projection head} that maps the original CLIP embeddings into the sparse concept space. The projection head produces a non-negative weight for every concept, yielding an interpretable embedding whose active dimensions indicate the semantic concepts most strongly associated with the image or text.


\subsubsection{Objective Function}

\paragraph{FILIP} (Fine-grained Interactive Language-Image Pre-training) introduces a fine-grained supervision mechanism. 
Instead of representing an image and a text by a single global embedding, FILIP preserves the embeddings of individual image patches 
and text tokens, enabling interactions at a finer semantic level.

After the image and text encoders produce patch-level and token-level embeddings, FILIP computes a pairwise similarity matrix between every image patch and every text token. 
Rather than directly comparing global representations as in CLIP, the similarity between an image and a text is obtained through a \emph{cross-modal late interaction} mechanism. 
Specifically, for each image patch, the maximum similarity over all text tokens is computed, while for each text token, the maximum similarity over all image patches is also calculated. 

The resulting similarities are then aggregated (using max pooling followed by averaging) to obtain the final image-text similarity score.

This objective encourages each image region to align with its most relevant textual token and vice versa, 
leading to more accurate word-patch correspondences.

\paragraph{}

